\documentclass[leqno, a4paper, 12pt]{jsarticle}
\usepackage{amssymb}
\usepackage{amsthm}
\usepackage{amsmath}
\usepackage{comment}
\usepackage{url}

\usepackage{enumitem}
\usepackage{showkeys}


%定理環境 thmはあまり使わずpropをなるべく使う。
%主張は基本的に証明の中で使い、そのため番号を付けない。
%注意環境を付けてみた。
\newtheorem{thm}{定理}
\newtheorem{prop}[thm]{命題}
\newtheorem{cor}[thm]{系}
\newtheorem{lem}[thm]{補題}
\newtheorem{df}[thm]{定義}
\newtheorem{exam}[thm]{例}
\newtheorem{fact}[thm]{事実}
\newtheorem*{claim}{主張}
\newtheorem*{cau}{注意}
\newtheorem*{rmk}{注意}
\renewcommand{\proofname}{証明}
%矢印たち。
%\toは\rightarrowと同じ。\getsは\leftarrowと同じ。同値は\iff
\newcommand{\To}{\Longrightarrow}
\newcommand{\Gets}{\LongLeftarrow}
%黒板太字
\newcommand{\nn}{\mathbb{N}}
\newcommand{\zz}{\mathbb{Z}}
\newcommand{\qq}{\mathbb{Q}}
\newcommand{\rr}{\mathbb{R}}
\newcommand{\cc}{\mathbb{C}}
%無限大
\newcommand{\mgn}{\infty}
%差集合は\setminusで統一する。等号の否定は\neq
%真部分集合は\subsetneqq　偏微分は\partial
%ディスプレイスタイル。\dfracでディスプレイ分数
\newcommand{\dpst}{\displaystyle}

\title{論文メモ
\large{\\--- エルデシュ的な話：普遍空間？ ---}
}
\author{ナンブ　キトラ}
\date{\today}
\begin{document}
\maketitle

本棚を整理したいので，
溜まっている論文のメモを書いて未来への手紙とすることにする．

普遍空間の話だと思う．

性質$(S_{n})$
というのは
ある$\eta_{n}$が存在して
$\rr$の$n$個の
点で直径が$\eta_{n}$より小さいものを与えたときにそれの
並行移動を含むという
性質である．

論文
\cite{MR0089886}
においてエルデシュと角谷は
次のような$\rr$
の部分集合$S$を構成した．
1. $S$は完全集合，
2. $\rr$の有限集合$F$を与えると，
適当な正の実数$r$と
実数$t\in \rr$
が存在して
$ F+t\subset S$
となる(任意の$n$について性質$(S_{n})$を満たす．)．
つまり，$\rr$
の全ての有限部分集合を拡大縮小並行移動を通して全て知っている完全集合が$S$
である．
詳しい話は
\cite{plewik2003set}
にも載っている．
論文
\cite{MR0112938}
もなんかこういう話で
性質$(S_{n})$
以外の並行移動に関する性質を
持つような実数の部分集合を構成している．

論文
\cite{MR2431353}
は逆に与えられたパターンを含まないようなハウスドルフ次元1の実数の
部分空間の構成である．

論文
\cite{MR0366673}
は何か色々参考文献があって面白いような気がするが何やってんのかよくわかんなかった．
多分ちゃんと読むといい気がする．
カントール集合が性質$S_{3}$
を満たさない理由を説明しているらしい．

論文
\cite{MR1417342}
ではシュタインハウス性質について研究してて
シュタインハウス性質をもつ集合はない点を持たないこととか
閉集合でないことを証明しているようだ．
ちょっと論文の書き方が訳わからんのだが，
平面の部分集合が
\textbf{シュタインハウス性質を持つ}
とは
多分，どんなに並行移動させたり回転させても
$\zz^{2}$との共通部分が必ず一点集合になることだと思うけど自信はない．

%READ!
%myplain is the same to amsplain style. 
\nocite{*}
\bibliographystyle{myplaindoidoi}
\bibliography{Eruniv.bib}



\end{document}